% =====================================================================
%  ch8_epilogue.tex — §8 结语：高维情形
% =====================================================================
\section{结语：高维情形}\label{s:higherdim}

我们预期本文对曲线及其镜像所述的构造与结果容许高维推广；不过写作时细节仍是尝试性的，以至于我们甚至不提出精确的猜想。

在 B 侧，考虑容许\emph{广义裤子分解}的刚性解析空间 $X_K$，即开覆盖由解析子集 $U_v$ 组成，而这些子集由 $n$ 维裤子
\[
\Pi_n=\bigl\{(z_0:z_1:\dots:z_{n+1})\in \PP^{n+1}\ \big|\ \textstyle\sum z_i=0,\ z_j\neq 0\ \forall j\bigr\}
\]
（即 $\PP^n$ 中 $n+2$ 个一般位置超平面之补）附加坐标赋值上的适当不等式得到。这类分解最常见地来自复簇到热带极限的极大退化。分解的组合结构由 $X_K$ 的\emph{热带化}编码，它是一个多面体复形 $\Sigma$：其顶点标记 $X_K$ 分解中的裤子 $U_v$，而高维胞腔及其仿射结构决定哪些子集 $U_v$ 有非平凡交叠，以及坐标赋值在黏合映射下如何变换。

每个裤子 $U_v$ 有 $\binom{n+2}{n}$ 个\emph{端} $U_{\sigma/v}$，即 $U_v$ 中使齐次坐标中有 $n$ 个小于其余两个的子集；每个端对应在 $v$ 处相交的 $\binom{n+2}{n}$ 个最高维胞腔 $\sigma\subset \Sigma$ 之一。又，$\Sigma$ 的每个最高维胞腔 $\sigma$ 决定子集 $U_\sigma\subset X_K$，沿它相邻于 $\sigma$ 的诸裤子 $U_v$ 的端 $U_{\sigma/v}$ 相互交叠。适当选取坐标后，$U_\sigma$ 可（在保值的方式下）与由 $\sigma$ 上仿射结构决定的仿射域 $\val^{-1}(\sigma)\subset (K^*)^n$ 等同。这又给出 $X_K$ 分解中每个裤子的每个端 $U_{\sigma/v}$ 上的 $(K^*)^n$-值局部坐标 $t_{\sigma/v}$；按构造，这些坐标的赋值与 $U_v$ 嵌入的模型裤子 $\Pi_n$ 上齐次坐标之比的赋值相同，且在相邻于给定最高维胞腔 $\sigma$ 的不同顶点 $v$ 之间按单项式坐标变换变化。（注意 $X_K$ 的完整描述还涉及沿 $\Sigma$ 中所有维数胞腔的黏合，为简单起见从略。）

在 A 侧，考虑由环面 K\"ahler 流形 $M_\sigma$ 沿环面层粘成的分层空间 $M$，其沿余维 $k$ 层的局部模型是 $(\C^*)^{n-k}$ 与 $\C^{k+2}$ 中所有坐标 $k$-平面之并 $\amalg_k$ 的乘积（特别地，每个顶点处有 $\binom{n+2}{n}$ 个最高维层相交）。为镜像对称之用，$M$ 各分量的矩映射像应与多面体复形 $\Sigma=\mathrm{Trop}(X_K)$ 的层及其仿射结构相符。

于是可考虑分层拉格朗日子流形 $L\subset M$：沿 $M$ 的每个余维 $k$ 层，它以 $(\C^*)^{n-k}$ 的光滑拉格朗日子流形与 $\amalg_k$ 中实正迹之并的乘积为模型。还可进一步要求 $L$ 的每个分量 $L_\sigma\subset M_\sigma$ 是矩映射纤维化 $M_\sigma\to \sigma$ 的截面，并给 $L$ 配上酉秩 $1$ 局部系统 $\cE$（在顶点处平凡化）。我们预期这类对象对应 $X_K$ 上的线丛：在每个裤子 $U_v$ 上取平凡线丛，再沿 $U_\sigma$ 用转移函数黏合，而转移函数由 $L$ 关于 $M$ 各层中实正迹 $L_0$ 的旋转数、$L$ 与 $L_0$ 围出的三角区域带符号辛面积、以及局部系统 $\cE$ 的和乐显式决定，正如第~\ref{ss:assignbundles} 节。

与曲线情形一样，预期 $M$ 的 Fukaya 范畴中态射空间的定义涉及 Hamilton 扰动，其流在每个低维层附近把拉格朗日的渐近方向旋转一个小的正量（当 $M$ 非紧时还涉及无穷远处的缠绕）。Fukaya 范畴的结构映射则应涉及（扰动）全纯圆盘刚性构形的加权计数，这些圆盘可经低维层在 $M$ 的各层之间传播。最简单、也是这里唯一讨论的情形，是 Floer 轨道从最高维层 $M_{\sigma_{in}}$ 经两层共有顶点 $v$ 传播到另一个最高维层 $M_{\sigma_{out}}$。此时全纯圆盘在 $v$ 附近的局部行为可这样描述：给入分量配置次数 $k_{in}\in \Z_{>0}^n$，给出分量配置次数 $k_{out}\in \Z_{\geq 0}^n$。我们进而预期传播系数 $C^{v;\sigma_{in},\sigma_{out}}_{k_{in},k_{out}}$（传播 Floer 轨道重数的局部贡献）应定义为单项式 $t^{k_{out}}_{\sigma_{out}/v}$ 在 $t_{\sigma_{in}/v}^{-k_{in}}$ 的（解析延拓）按 $U_{\sigma_{out}/v}$ 上坐标 $t_{\sigma_{out}/v}$ 的幂级数展开中的系数。

% [Tutor Note: 直觉] 本节只给「怎么猜」，未给定理：分层拉格朗日、转移函数、传播系数的定义都是按曲线情形逐字照搬的，但高维时圆盘可沿任意维数的层传播，局部模型也从 $(\C^3,-xyz)$ 变成 $(\C^*)^{n-k}\times\amalg_k$，$\mu^k$ 的 $A_\infty$ 关系是否成立并未验证。读者应把本节当作研究纲领，而不是已证结论。依据：原文 §8 末段（本工程 01_通读笔记）。

显然，还需大量工作来充实此故事的关键细节，并确认所提构造可靠且导致同调镜像对称的陈述；这项工作仍处于早期阶段，但希望读者已被说服：本文展开的故事很可能超出曲线情形。